\section{Conservation Laws}

\label{sec:conservation}

\subsection{Energy Conservation (No-Arbitrage)}

\begin{theorem}[Energy Conservation]
\label{thm:energy}
Along any Hamiltonian trajectory, $\mathcal{H}(\bm{q}(t), \bm{p}(t)) = \kappa$ for all $t$, where $\kappa$ is determined by initial conditions.
\end{theorem}

\begin{proof}
By direct computation:
\begin{align}
\frac{d\mathcal{H}}{dt} &= \sum_i \left(\frac{\partial \mathcal{H}}{\partial q_i}\dot{q}_i + \frac{\partial \mathcal{H}}{\partial p_i}\dot{p}_i\right) \\
&= \sum_i \left(\frac{\partial \mathcal{H}}{\partial q_i}\frac{\partial \mathcal{H}}{\partial p_i} - \frac{\partial \mathcal{H}}{\partial p_i}\frac{\partial \mathcal{H}}{\partial q_i}\right) = 0.
\end{align}
\end{proof}

\begin{corollary}[No-Arbitrage]
For any sequence of swaps $\{(\Delta\bm{q}^{(k)}, \Delta\bm{p}^{(k)})\}_{k=1}^K$ that returns the pool to its initial state, the net credit transfer satisfies $\sum_k \Delta\bm{p}^{(k)} \cdot \bm{f}^{(k)} > 0$, where $\bm{f}^{(k)}$ are the fee vectors.
\end{corollary}

This follows because energy conservation prevents ``free'' round-trip extraction, and the addition of positive fees makes any cycle strictly costly.

\subsection{Liouville's Theorem (Volume Preservation)}

\begin{theorem}[Phase Space Volume Preservation]
\label{thm:liouville}
The Hamiltonian flow $\phi_t : \Gamma \to \Gamma$ preserves the phase space volume form $\Omega = \omega^n / n!$:
\begin{equation}
\phi_t^* \Omega = \Omega \quad \forall t.
\end{equation}
\end{theorem}

\begin{proof}
The divergence of the Hamiltonian vector field $X_\mathcal{H} = (\partial_p \mathcal{H}, -\partial_q \mathcal{H})$ vanishes:
$\nabla \cdot X_\mathcal{H} = \sum_i (\partial^2 \mathcal{H}/\partial q_i \partial p_i - \partial^2 \mathcal{H}/\partial p_i \partial q_i) = 0$.
By the Liouville equation, $d\Omega/dt = -(\nabla \cdot X_\mathcal{H})\Omega = 0$.
\end{proof}

\textbf{Economic interpretation}: Volume preservation means the pricing function is \emph{information-preserving}. No market state information is destroyed by the dynamics, ensuring that historical price data can be used for accurate statistical inference.

\subsection{Angular Momentum (Reserve Balance)}

For the symmetric case ($w_i = 1/n$ and $\bar{q}_i = \bar{q}$ for all $i$), define the angular momentum:
\begin{equation}
L_{ij} = q_i p_j - q_j p_i.
\end{equation}

\begin{proposition}
If $w_i = w_j$ and $\bar{q}_i = \bar{q}_j$, then $\{L_{ij}, \mathcal{H}\} = 0$, so $L_{ij}$ is conserved.
\end{proposition}

\textbf{Economic interpretation}: When two resources have identical weights and targets, the relative imbalance between them is conserved. This prevents one resource from systematically draining the other through the AMM mechanism.

\subsection{Noether's Theorem and Economic Symmetries}

By Noether's theorem, every continuous symmetry of the Hamiltonian corresponds to a conserved quantity. We catalog the relevant symmetries:

\begin{table}[H]
\centering
\begin{tabular}{lll}
\toprule
\textbf{Symmetry} & \textbf{Conserved Quantity} & \textbf{Economic Meaning} \\
\midrule
Time translation & $\mathcal{H}$ & No-arbitrage \\
Rotation (sym.) & $L_{ij}$ & Reserve balance \\
Scaling & Total value & Market cap invariance \\
Phase space vol. & $\Omega$ & Information preservation \\
\bottomrule
\end{tabular}
\caption{Noether symmetries and their economic interpretations.}
\label{tab:noether}
\end{table}

